\chapter{Óptimos de planeo, viento y altitud}\label{ch:planeo-punto}
\objetivos{Derivar las incidencias de máximo alcance y mínimo descenso; imponer restricciones; interpretar la modificación del óptimo por viento.}

\section{Mejor planeo en calma}
Con polar $C_D=C_{D0}+kC_L^2$, la eficiencia es
\begin{equation}
E_a(C_L)=\frac{C_L}{C_{D0}+kC_L^2}.
\end{equation}
Su derivada vale
\begin{equation}
\frac{\dd E_a}{\dd C_L}=
\frac{C_{D0}-kC_L^2}{(C_{D0}+kC_L^2)^2}.
\end{equation}
Para $C_L>0$, crece antes de $\sqrt{C_{D0}/k}$ y decrece después. El máximo interior es
\begin{equation}
C_{L,E}=\sqrt{C_{D0}/k},\quad C_{D,E}=2C_{D0},\quad
E_{a,\max}=\frac1{2\sqrt{C_{D0}k}}.\label{eq:ldmax}
\end{equation}
La igualdad de las dos contribuciones $C_{D0}=kC_L^2$ se refiere a los dos términos de la polar total adoptada. Sólo es igualdad de resistencia parásita e inducida pura si el término cuadrático contiene únicamente esta última.

El ángulo de descenso mínimo exacto es $\delta_{\min}=\arctan(1/E_{a,\max})$. El $C_L$ que maximiza el alcance en calma es el mismo en la formulación exacta y en la de ángulo pequeño, porque ambas utilizan la misma razón $L/D$.

\section{Mínima tasa de descenso}
En la aproximación tendida se minimiza
\begin{equation}
f(C_L)=\frac{C_{D0}+kC_L^2}{C_L^{3/2}}
=C_{D0}C_L^{-3/2}+kC_L^{1/2}.
\end{equation}
La condición $f'=0$ conduce a $kC_L^2=3C_{D0}$. Por tanto,
\begin{equation}
C_{L,s}=\sqrt{3C_{D0}/k}=\sqrt3C_{L,E},\qquad
V_s\simeq3^{-1/4}V_E.
\label{eq:min-sink}
\end{equation}
Aquí $V_s$ significa velocidad de mínimo descenso sólo en esta ecuación; para evitar confusión, en tablas se escribirá $V_{\rm min\,desc}$ y se reservará $V_{\rm stall}$ para pérdida. La eficiencia aerodinámica en mínimo descenso es $\sqrt3/2$ del máximo, por lo que el planeo resulta algo más inclinado aunque la velocidad vertical sea menor.

\subsection{Corrección exacta}
Si se minimiza la expresión exacta del capítulo anterior,
\begin{equation}
\frac{2kC_L}{C_D}-\frac{3}{2}
\frac{C_L+2kC_LC_D}{C_L^2+C_D^2}=0.
\end{equation}
Para $C_L>0$ se reduce a
\begin{equation}
4kC_L^2-2kC_D^2-3C_D=0.
\label{eq:min-sink-exacto}
\end{equation}
Al despreciar el término de orden $C_D^2$ y sustituir la polar, se recupera $kC_L^2\simeq3C_{D0}$. Se debe buscar el mínimo dentro del intervalo válido de la polar: una raíz lejana asociada a sustentación enorme no representa un planeador operativo. Esta distinción evita presentar una aproximación excelente para planeo tendido como identidad universal.

\section{Velocidades equivalentes y restricciones}
Las velocidades TAS de los óptimos son proporcionales a $\sqrt{W/\rho}$. Sus EAS son aproximadamente proporcionales a $\sqrt W$. Si el óptimo requiere $C_L>C_{L\max}$, el máximo factible aparece en la frontera de sustentación, con el margen operativo exigido. No basta con imprimir la raíz de la derivada.

\caso{Ejemplo original: comparación de dos objetivos}{Con la polar del capítulo anterior, $C_{L,E}=\sqrt{0.013/0.031}=0.6476$ y $C_{L,s}\simeq1.1216$. Ambos caben en un modelo con $C_{L\max}=1.6$. El primer valor maximiza distancia desde una altura en calma; el segundo minimiza descenso en la aproximación tendida. La relación de velocidades es aproximadamente $0.7598$. Si una configuración tuviese $C_{L\max}=0.95$, el mínimo descenso interior dejaría de ser alcanzable.}

\ilustra{eficiencias}{Tres funciones objetivo para una misma polar. Sus máximos difieren: eficiencia, mínimo consumo de potencia y alcance específico de un reactor bajo las hipótesis que se desarrollan después.}

\section{Mejor alcance sobre el suelo con viento}
Supóngase viento horizontal alineado con la ruta, uniforme y constante. Se escribe $U>0$ para viento de cola. Con una tasa de descenso positiva $s(V)$ y planeo tendido, el avance por altura es
\begin{equation}
G(V)=\frac{V+U}{s(V)}.
\end{equation}
La condición interior es
\begin{equation}
s(V)-(V+U)s'(V)=0.\label{eq:optimo-viento}
\end{equation}
En calma se recupera la tangencia desde el origen a la polar de descenso. Con viento de cara $U<0$, el origen de la tangencia se desplaza hacia velocidad positiva y el óptimo suele aumentar: se vuela más rápido para reducir el tiempo durante el que el viento resta avance. Con viento de cola el óptimo suele disminuir, dentro de los límites de pérdida y control.

La versión exacta sustituye $V$ por $V_H(V)$ y deriva $[V_H(V)+U]/s(V)$. Si el parámetro de la curva es $C_L$, se deriva respecto a él. El cambio de variable no altera el óptimo si la correspondencia es regular en el intervalo considerado.

\section{Viento vertical y autonomía}
En una masa de aire con velocidad vertical ascendente $U_h$, la variación de altitud sobre el suelo es $\dot h=U_h-s$. Un planeador puede mantener o ganar altura si $U_h\geq s$. Ya no procede calcular una autonomía finita como $\Delta h/s$ usando sólo la tasa de descenso en aire en calma. Además, un campo de ascendencias no suele ser uniforme ni indefinido; el objetivo incluye buscar y permanecer en regiones de energía atmosférica.

El mínimo descenso en el aire no cambia por añadir una velocidad vertical uniforme independiente de la velocidad del avión. Sí cambia la autonomía geométrica sobre el terreno y puede cambiar la estrategia si las ascendencias ocupan regiones limitadas. En esta edición los cálculos integrales de planeo básico suponen viento vertical nulo.

\section{Viraje de un planeador}
En un descenso helicoidal tendido con alabeo aproximadamente constante, la sustentación requerida aumenta aproximadamente a $nW$, con $n\simeq1/\cos\phi$. A igual $C_L$, la velocidad aumenta como $\sqrt n$ y la potencia resistente como $n^{3/2}$ respecto al caso rectilíneo del mismo peso, bajo la polar y aproximaciones adoptadas. Por tanto, permanecer girando suele aumentar el descenso propio del avión. El beneficio dentro de una térmica depende de la distribución de velocidad ascendente y del radio de giro, no sólo de minimizar la resistencia.

Para descensos grandes deben conservarse $\cos\gamma$ en el equilibrio vertical y las relaciones de curvatura de la trayectoria. La fórmula nivelada $n=1/\cos\phi$ no se presenta como exacta para cualquier hélice descendente.

\section{Taller resuelto: Óptimo exacto de mínimo descenso}

\noindent\textbf{Nivel:} Intermedio.

\apartado{1. Enunciado}{Con el planeador anterior, encuentre por cálculo numérico el CL de mínimo descenso exacto en 0.3<=CL<=1.55 y compare con sqrt(3CD0/k). Compruebe la condición de derivada exacta.}

\apartado{2. Datos}{Masa 455 kg y densidad 1.05 kg/m³; mismo dominio válido de polar.}

\apartado{3. Incógnitas}{Las magnitudes y comprobaciones solicitadas en el caso.}

\apartado{4. Sistema de referencia}{Planeo rectilíneo en calma; ángulo de descenso positivo delta.}

\apartado{5. Hipótesis físicas}{Equilibrio exacto y polar fija; se optimiza sólo dentro de su intervalo.}

\apartado{6. Modelo matemático}{Óptimo exacto de mínimo descenso}

\apartado{7. Ecuaciones simbólicas}{Se parte de los balances y definiciones siguientes.}

\begin{equation}s=\sqrt{2W/(\rho S)}\,C_D/(C_L^2+C_D^2)^{3/4}\end{equation}

\begin{equation}4kC_L^2-2kC_D^2-3C_D=0\end{equation}

\apartado{8. Desarrollo algebraico}{Las operaciones se realizan en este orden, conservando las condiciones del modelo.}

\begin{enumerate}

\item Se forma la función exacta de descenso sin aproximar coseno.

\item Se busca el mínimo acotado y se comparan extremos.

\item Se evalúa el residuo de la condición algebraica exacta cerca del óptimo.

\item Se compara con el CL de la fórmula tendida y se cuantifica la diferencia.

\end{enumerate}

\apartado{9. Sustitución numérica}{Se usan las magnitudes del apartado de datos en SI; los ángulos se convierten a radianes antes de evaluar funciones trigonométricas. Las densidades indicadas como ISA se obtienen de la capa correspondiente.}

\noindent CL exacto:\quad $1.148396272098129\simeq 1.148396\,$.\par

\noindent CL aproximado:\quad $1.14564392373896\simeq 1.145644\,$.\par

\noindent Descenso mínimo exacto:\quad $1.044153049097051\simeq 1.044153\,\mathrm{m\,s^{-1}}$.\par

\apartado{10. Resultados con unidades}{Los resultados pertenecen al modelo y al intervalo de datos declarados.}

\begin{center}\begin{tabular}{lr}\toprule Magnitud&Resultado\\\midrule

CL exacto & $1.148396\,$\\

CL aproximado & $1.145644\,$\\

Descenso mínimo exacto & $1.044153\,\mathrm{m\,s^{-1}}$\\

Diferencia relativa CL & $0.002402447\,$\\

\bottomrule\end{tabular}\end{center}

\apartado{11. Interpretación}{El óptimo de la aproximación y el exacto se aproximan para ángulos pequeños; no se los identifica sin evaluar el error.}

\apartado{12. Comprobación independiente y alcance}{Búsqueda de raíz de 4kCL²-2kCD²-3CD y comparación con minimización directa de la tasa.}

\apartado{13. Limitaciones}{Una segunda raíz fuera del dominio no representa una condición de vuelo admisible.}

\par\smallskip\noindent{\footnotesize\textit{Datos y solución originales. Las sustituciones registradas se reproducen con SymPy y Studium; eso no establece validación experimental ni revisión experta.}}


\repaso{¿Por qué no basta con maximizar $L/D$ para regresar contra viento?}{Porque el alcance sobre el suelo incluye el tiempo de exposición al viento. La razón relevante es avance terrestre dividido por descenso; la condición de tangencia cambia.}
\resumen{Mejor planeo y mínimo descenso son problemas de optimización diferentes. El viento modifica el objetivo de alcance terrestre. Las raíces interiores se comparan con pérdida, límites de polar y restricciones de misión antes de interpretarlas como velocidades utilizables.}
\fuentes{How, MIT 16.333, lección 1, actuaciones de planeo. La condición con viento, la corrección exacta y los ejemplos se derivan en esta edición. La clasificación de ejercicios se basa en los cuadernos facilitados, sin adoptar sus soluciones como autoridad.}
